\chapter*{Заключение}
\addcontentsline{toc}{chapter}{Заключение}
\markboth{Заключение}{}

Двенадцать глав назад мы начали с вопроса, что происходит с программой между
нажатием кнопки «запустить» и появлением результата на экране. Теперь на него есть
развёрнутый ответ: исходный текст компилируется в байт-код, загрузчик классов
подтягивает нужные классы в память виртуальной машины, интерпретатор начинает
выполнение, а JIT-компилятор постепенно переводит горячие участки в машинный код.
Понимание этой цепочки~--- не абстрактное знание: именно оно объясняет, почему
Java-приложение медленно стартует и быстро работает потом, почему утечка памяти
возможна и в языке со сборщиком мусора, и почему один и тот же код ведёт себя
одинаково на разных операционных системах.

Но главное, чему стоило научиться,~--- это не синтаксис. Синтаксис забывается и
гуглится. Ценность представляют привычки, которые формируются по ходу:
проверять граничные случаи, а не только счастливый путь; отделять данные от
логики, а логику~--- от представления; писать код так, чтобы через полгода его
понял другой человек (чаще всего этим человеком оказываетесь вы сами);
не оптимизировать то, что не измерено; и коммитить достаточно часто, чтобы
никогда не бояться сломать работающее.

Отдельно стоит сказать о том, чего в этой книге нет. Мы почти не касались
микросервисной архитектуры, контейнеризации, реактивного программирования,
облачных развёртываний и производительности под нагрузкой~--- каждая из этих тем
заслуживает собственного пособия. Не разобраны и целые пласты экосистемы: Kotlin
как альтернативный язык для той же виртуальной машины, Gradle в его продвинутых
сценариях, Kafka и очереди сообщений, GraalVM и компиляция в native-образ. Это
не пробелы, а естественная граница вводного курса. Хорошая новость в том, что
фундамент, заложенный здесь, делает все эти темы доступными: они надстраиваются
над тем, что уже освоено, а не заменяют его.

Что делать дальше? Самый честный ответ~--- писать код. Учебные задания отличаются
от настоящих тем, что в них заранее известно, что требуется сделать. В реальном
проекте половина работы уходит на то, чтобы понять задачу. Поэтому лучшее
продолжение этого курса~--- собственный проект: что-нибудь небольшое, но доведённое
до состояния, когда им пользуется кто-то кроме автора. Такой проект за месяц даёт
больше, чем ещё один прочитанный учебник.

Если же хочется системного продолжения, разумный порядок такой: сначала углубить
Spring (безопасность, кеширование, работа с очередями), затем разобраться
с контейнерами и развёртыванием, и только потом браться за распределённые
системы. И на каждом шаге~--- возвращаться к основам из первых глав. Опыт
показывает, что большинство трудных ошибок в сложных системах объясняется не
экзотикой фреймворков, а забытыми деталями работы виртуальной машины,
многопоточности и транзакций.

Успехов в разработке.
